\documentclass[12pt,a4paper]{article}

\usepackage{cmap}
\usepackage[T2A]{fontenc}
\usepackage[utf8]{inputenc}
\usepackage[russian]{babel}
\usepackage{amsmath,amssymb}
\usepackage[a4paper,margin=2cm]{geometry}

\setlength{\parindent}{0pt}
\setlength{\parskip}{0.45em}
\sloppy

\begin{document}

\begin{center}
{\Large\bfseries Билет 14: вопросы для проверки знаний}
\end{center}

\noindent\fbox{%
  \begin{minipage}{\dimexpr\linewidth-2\fboxsep-2\fboxrule\relax}
  \normalfont
Экстремумы функций многих переменных: необходимое условие, достаточное
условие.
  \end{minipage}%
}

\section*{Вопросы на понимание}

\begin{enumerate}
\item Чем локальный минимум функции на множестве \(X\) отличается от безусловного локального минимума?
\item Почему точка на границе области может быть точкой локального экстремума, хотя необходимое условие \(\operatorname{grad} f=0\) к ней напрямую не применяется?
\item Чем строгий локальный экстремум отличается от нестрогого? Приведите пример нестрогого минимума, который не является строгим.
\item Почему условие \(\operatorname{grad} f(x^0)=0\) является только необходимым, а не достаточным для экстремума?
\item Может ли функция иметь локальный экстремум в точке, где она не дифференцируема? Как это связано с необходимым условием?
\item Как рассмотрение сечений \(f(x^0+t e_i)\) помогает понять равенство нулю всех частных производных в точке безусловного экстремума?
\item Почему в достаточном условии нужна стационарность точки \(x^0\), а не только знак второго дифференциала?
\item Что означает положительная определенность второго дифференциала \(d^2f(x^0)[h]\) с точки зрения приращений \(f(x^0+h)-f(x^0)\)?
\item Для функций \(f(x,y)=x^2+y^2\), \(g(x,y)=-x^2-y^2\) и \(u(x,y)=x^2-y^2\) в точке \((0,0)\) какой тип квадратичной формы второго дифференциала возникает и какой вывод об экстремуме ожидается?
\item Почему знаконеопределенная квадратичная форма исключает локальный минимум и локальный максимум?
\item Что может пойти не так в вырожденном случае, когда второй дифференциал полуопределен или равен нулю? Сравните поведение \(x^4\) и \(x^3\) в нуле.
\item Как изменится проверка экстремума, если ограничиться не всеми направлениями \(h\), а только несколькими выбранными прямыми?
\end{enumerate}

\section*{Источники для сверки}

\texttt{target/answer\_question\_14.tex};
\texttt{IvGE\_2.pdf}, стр. 45--48;
\texttt{IvGE\_1.pdf}, стр. 105--108.

\end{document}
